\begin{lemma}[Interchangeability at a positive hierarchy stage]
Let $T$ be a positive-rung ETT of rung number $n$, constructed from an
arbitrary certified input coloring and with generating coloring
$\varphi_n$.  Assume that $T$ has maximum property, that $\varphi_n$ is
proper and leaves exactly the distinguished edge $e_0$ uncolored, and
fix a hierarchy base.  Let
\[
 T_n=T_{n,0}\subset T_{n,1}\subset\cdots\subset
 T_{n,q}\subset T_{n,q+1}=T
\]
be a good hierarchy $H$ over that base, with its actual family of
two-color sets $\Gamma^j_\eta$.

Assume, uniformly over every certified input coloring and generating
coloring, that every positive-rung ETT $U$ of rung number $s<n$ with
maximum property satisfies all six lower-rung assertions:
\begin{enumerate}
\item $V(U)$ is elementary under its generating coloring;
\item under every proper coloring leaving exactly $e_0$ uncolored and
stable on the final ladder closure, that closure is interchangeable for
two colors whenever at least one is missing there;
\item stability descends at every PE support to the prefix before that
support and to its first-appearance ladder stage, and the stored pairs at
distinct PE slots with distinct supports and the same first-appearance
index are disjoint;
\item at a terminal PE, the stored two-color component through the
support meets the terminal ladder tree only at the support;
\item every nonminimum endpoint of a defective boundary color precedes
each later PE support or SE extension vertex allowed by its
first-appearance index; and
\item every proper terminal stable coloring leaving exactly $e_0$
uncolored is a mod-coloring of $U$.
\end{enumerate}

Assume also the induction hypothesis (6.1) in the following explicit
form.  For every arbitrary certified input coloring and generating
coloring $\psi$, every ETT $U$ of the same rung number $n$ as $T$ which
has maximum property, and every good hierarchy $K$ of $U$ with strictly
fewer than $q$ dividers, $V(U)$ is elementary under $\psi$.  Finally
assume the large-branch estimate
\[
              |\overline{\varphi_n}(T_n)|\geq 2n+11.       \tag{L}
\]

Let $1\leq p\leq q$, let
$S=T_{n,p}$ (formally, for $i=p-1$ this is
`$H.\mathit{stage}(i.\mathit{succ}.\mathit{castSucc})$'), and let
$\alpha,\beta\in[k]$ be distinct.  If
$\alpha\in\overline{\varphi_n}(S)$ and the singleton
$\{\alpha\}$ is closed on $S$ under $\varphi_n$, then $\alpha$ and
$\beta$ are $S$-interchangeable under $\varphi_n$.  All tree sequences,
paths, and recolorings retain the individual identities of parallel
edges.
\end{lemma}
\begin{proof}
Write $\varphi=\varphi_n$.  We first isolate the presentations on which
the induction will range.  If $1\leq r\leq q$, restrict the ordered
stage data of $H$ through $T_{n,r}$ and declare $T_{n,r}$ to be the
terminal tree.  The fields of \noderef{TashkinovHierarchy} restrict
entry by entry: the retained stages are
\[
 T_{n,0},\ldots,T_{n,r-1},T_{n,r},
\]
the retained divider positions are the first $r-1$, and the terminal
field becomes $T_{n,r}$.  Its divider count is therefore $r-1$.
The same certified series gives an \noderef{ExtendedTashkinovTree} on
$T_{n,r}$: the base and strict-prefix clauses follow from the base-stage,
core-inside-first, nested, and strict fields of
\noderef{TashkinovHierarchy}, while the upper initial-segment clause is
obtained by composing those ordered segments with the old ETT upper
segment.  The definition \noderef{MaximumProperty} depends only on the
certified ladder series, not on which intermediate terminal prefix is
displayed, so the given maximum-property certificate is unchanged.

Restrict the family $\Gamma$ to stages $0,\ldots,r-1$.  Pair sizes,
avoidance, and pairwise disjointness are literally the corresponding
fields of \noderef{GoodHierarchy}.  Its unique-refresh and minus-closed
fields are restricted to indices below $r-1$; its parallel-base field is
unchanged.  Thus this is a good hierarchy with $r-1<q$ dividers.  In
addition, the old `minusClosed' field at index $r-1$, after unfolding
\noderef{HierarchyStoredColorsAt} and terminal missing colors, says that
the new terminal tree is
$(\bigcup_{\eta\in D_{n,r}}\Gamma^{r-1}_\eta)^-$-closed.  We retain this
as a separate terminal-closure certificate; it is not incorrectly
claimed to be one of the good hierarchy's remaining internal fields.
The same-rung induction hypothesis (6.1), applied to this exact
truncated ETT and hierarchy, gives
\[
             V(T_{n,r})\text{ is elementary under }\varphi.
                                                               \tag{1}
\]
This construction works for arbitrary certified input and generating
colorings, which is essential when we repeat it after a recoloring.

We shall also truncate at a vertex $b$ lying strictly after
$A=T_{n,r-1}^{*}$ and no later than $T_{n,r}$.  Use
\noderef{TreeSequencePrefix} to take the ordered prefix through $b$, or
the prefix $B$ immediately before $b$.  Every retained vertex and edge
is the same indexed identity by \noderef{TreeSequenceInitialSegment}.
There are three genuinely different cases.

If $A$ is a strict initial segment of $B$, keep the stages through
$T_{n,r-1}$ and use $B$ as terminal.  The last strictness field of
\noderef{TashkinovHierarchy} is now exactly $A\subsetneq B$, so this is
the valid $r-1$-divider presentation.  If instead $B=A$ and $r>1$, do
not use that presentation: retain only the stages through
$T_{n,r-2}$ and take $T_{n,r-1}=A$ as terminal.  This is the valid
$r-2$-divider presentation, with its terminal minus-closed certificate
given by the old field at $r-2$.  In either case a recoloring on a full
component disjoint from $B$ is an \noderef{OutsideKempeRecoloring} of
this displayed ETT.  Once its two colors have been proved distinct,
\noderef{HierarchyLEMStable} supplies properness, exact
one-uncoloredness, the corresponding ETT, maximum property, the same
ordered stages, stored sets and $\Gamma$-pairs, and the transported
terminal minus-closed certificate.

The remaining case is $r=1$ and $B=A=T_{n,0}^{*}$.  We do not present
the starred base as an ETT.  The component is disjoint from $A$ and hence
from $T_n$.  Thus \noderef{OutsideKempeRecoloringStable} supplies the
proper, one-uncolored, terminal-ladder stability data.  Lower assertion
(vi), instantiated by \noderef{Td}, makes the changed coloring a
\noderef{ModColoring}; the zero-rung stability and interchangeability
endpoint is \noderef{TechnicalTheoremBase}.  This is the separate base
route used below.

For the one-edge enlargement from $B$ to the prefix $C$ through $b$,
the old final certified closure in
\noderef{TashkinovSeriesStep} and \noderef{TashkinovClosure} supplies the
exact finite augmentation sequence.  Because the recoloring fixes every
edge identity and every missing-color witness on $B$,
each retained augmentation still satisfies
\noderef{TashkinovAugmentation}; the universal clause of
\noderef{ModColoring}, obtained from lower assertion (vi) via
\noderef{Td}, then supplies the corresponding ETT on the enlarged
prefix.  \noderef{MaximumPropertyTransport} transfers maximum property.
When $B=A$ and $r>1$, adjoining $b$ restores the $r-1$-divider
presentation: the last strict containment is the exact one-edge
containment $A\subsetneq C$.  Pair size, disjointness, unique refresh,
parallel base, and every old minus-closed field are copied; final
avoidance is checked on that exact new edge.  When $r=1$, the same check
constructs the zero-divider hierarchy with base stage $T_n$, terminal
$C$, and core inside $C$.  Thus no hierarchy is ever asserted with its
last retained stage equal to its terminal tree, and no endpoint pair is
substituted for a graph-edge identity.

We now make the well-founded comparison domain explicit.  An
\emph{admissible counterexample presentation} consists of an arbitrary
certified input coloring, a proper generating coloring $\psi$ leaving
exactly $e_0$ uncolored, a same-rung ETT $U$ with maximum property, a
hierarchy base and good hierarchy $K$ with $d<q$ dividers, the same fixed
graph, palette, distinguished edge, lower six-assertion package and
(6.1), and the analogue of (L) for its own terminal ladder tree.  It also
contains two distinct colors $a,b$, with $a$ missing and singleton-closed
on $U$, a failure of $U$-interchangeability, and the separate terminal
minus-closed certificate
\[
 U\text{ is }
 \left(\bigcup_{\eta\in D^K_{n,d+1}}\Gamma^d_\eta\right)^-
 \text{-closed under }\psi.                                  \tag{C}
\]
Here $D^K_{n,d+1}$ is the terminal stored-color expression: unfolding
\noderef{TashkinovStoredColors} and
\noderef{HierarchyStoredColorsAt}, it is the cumulative raw stored
colors through rung $n$ with the colors missing on $U$ deleted.

Give such a presentation primary measure $d$.  Its secondary measure is
zero if $b$ is missing on $U$.  Otherwise the endpoint-parity argument
given in Case 2 below supplies three exit paths, and the secondary
measure is one plus the least possible sum of their lengths, where the
triple must have distinct outside endpoints and must lie in at least two
full two-color components.  The relevant sets of natural numbers are
nonempty, so well-ordering of $\mathbb N\times\mathbb N$ in lexicographic
order gives a least admissible presentation if any exists.

Suppose the stated conclusion is false.  The truncation through the
stated $T_{n,p}$, together with its old minus-closed field, is an
admissible presentation, so choose a lexicographically least one.  Rename
its generating coloring $\varphi$, its terminal tree $S$, its colors
$\alpha,\beta$, put $p=d+1$, and write
$A=T_{n,p-1}^{*}$ for its last starred stage.  All later recolored
presentations are compared in this same domain; the input coloring,
ETT, and hierarchy witnesses are allowed to change, whereas
$G,k,e_0,n$, the lower package, (6.1), $q$, and the large estimate are
the fixed parameters.

By (1), $S$ is elementary.  Since $\alpha$ is missing on $S$, it is
missing at exactly one vertex: existence is
\noderef{TreeMissingColors}, and uniqueness is
\noderef{ColoredVertexSetElementary}.  Properness says that every other
vertex is incident with exactly one $\alpha$-edge.  The singleton
closedness hypothesis, expanded using \noderef{TreeColorSetClosed} and
\noderef{BoundaryColorEdges}, says that no such edge is a boundary edge.
Counting incidences of exact $\alpha$-edge identities inside $S$ therefore
gives
\[
                         |V(S)|\text{ is odd}.                \tag{2}
\]
Indeed, the $|V(S)|-1$ present incidences are paired by the two endpoints
of each internal $\alpha$-edge.  Finally terminal certificate (C) of the
chosen admissible presentation gives
\[
 S\text{ is }
 \left(\bigcup_{\eta\in D_{n,p}}\Gamma^{p-1}_\eta\right)^{-}
 \text{-closed under }\varphi.                               \tag{3}
\]
This is exactly the finite-set predicate
\noderef{TreeColorSetMinusClosed}, not an unstated closure convention.

The negation of \noderef{TreeInterchangeable} supplies two distinct
nontrivial $(\alpha,\beta)$ path components meeting $S$.  By
\noderef{TwoColorPathComponent} each has a genuine edge node and a path
endpoint; by \noderef{TwoColorComponentSpanningPath} it has a simple
spanning path, and \noderef{TwoColorIncidenceDegreeAtMostTwo} controls
all its incidences.  Since $\alpha$ has no boundary edge, every crossing
from $S$ is an exact $\beta$-colored edge identity.  The endpoint and
parity counts in the two cases below show how many outside endpoint
portions occur; each such portion ends at its first vertex of $S$, so it
satisfies \noderef{TreeExitPath}.  Distinct components are vertex-disjoint
by \noderef{TwoColorComponentsDisjoint}.  These observations justify all
exit paths below directly in the incidence graph
\noderef{TwoColorIncidenceGraph}, where parallel edges are separate
nodes.

\textbf{Case 1: $\beta\in\overline\varphi(S)$.}
Elementarity makes $\beta$ missing at a unique vertex.  We do not assert
that this vertex differs from the unique $\alpha$-missing vertex: one
vertex may miss both distinct colors.  Counting $\beta$-incidences as
above nevertheless gives
\[
                  |\partial_{\varphi,\beta}(S)|\equiv
                  |V(S)|-1\equiv0\pmod2.                    \tag{4}
\]
Use \noderef{TwoColorComponentSpanningPath} on each of the two distinct
full witness components.  Properness and
\noderef{TwoColorIncidenceDegreeAtMostTwo} make each support a simple
path whose two end vertex-nodes miss $\alpha$ or $\beta$; edge-nodes have
incidence degree two and cannot be ends.  The two components are disjoint
by \noderef{TwoColorComponentsDisjoint}, so their four endpoints are
four distinct vertices.  Elementarity permits at most one of them in
$S$ missing $\alpha$ and at most one in $S$ missing $\beta$.  These two
vertices may coincide, but in either event at least two of the four
endpoints lie outside $S$.

From two outside endpoints follow their spanning paths toward their first
vertex-nodes in $S$.  Such a first contact exists because each selected
full component meets $S$.  The resulting portions satisfy every clause
of \noderef{TreeExitPath}.  If the endpoints came from distinct full
components, their portions are disjoint by component disjointness.  If
they are the two ends of one component, they are the two end portions of
one simple spanning path.  Their first contacts cannot coincide: both
last edges would then be distinct $\beta$-colored edges incident with the
same vertex, since $\alpha$ has no boundary edge, contradicting
properness.  The two end portions are therefore vertex-disjoint also in
this case.  Call them $P_1,P_2$, write $a_j$ for the outside end and
$b_j$ for the inside end, and order them so that $b_1\prec b_2$ in
\noderef{TreeVertexOrder}.  We retain, for later use, whether these tails
came from two original full components or from the two ends of one.

\emph{Subcase 1.1: $b_2\notin V(A)$.}
The first edge of $P_2$ at $b_2$ is a $\beta$-colored boundary identity.
If $\beta$ were outside
$\bigcup_{\eta\in D_{n,p}}\Gamma^{p-1}_\eta$, (3) would make its boundary
empty.  Hence there is $\eta\in D_{n,p}$ with
$\beta\in\Gamma^{p-1}_\eta$.  Unfold
\noderef{TashkinovStoredColors} and
\noderef{HierarchyStoredColorsAt}.  Since $A\subseteq S$, every color
missing on $A$ is missing on $S$, so the terminal difference satisfies
$D_{n,p}\subseteq D_{n,p-1}$.  Thus $\eta$ is in the preceding stored
set, and \noderef{GoodHierarchyReservedDisjoint} puts this exact pair in
the colors missing on $A$.  Hence
\[
                         \beta\in\overline\varphi(A).         \tag{5}
\]

Choose $\gamma\in\overline\varphi(b_2)$.  This choice is justified from
the certificate rather than assumed: the initial constructor in
\noderef{TashkinovSeriesStep} contains a \noderef{KTriple}, hence
$k\geq\Delta(G)+1$; properness gives at most
$d_G(b_2)\leq\Delta(G)$ distinct present colors by \noderef{Degree} and
\noderef{MaximumDegree}; the complement definition
\noderef{MissingColors} is therefore nonempty.  If $\gamma$ belonged to
the reserved union in (3), the preceding-stage conclusion of
\noderef{GoodHierarchyReservedDisjoint} would make it missing at a
vertex of $A$.  That vertex differs from $b_2$, contradicting (1).
Consequently $\gamma$ is outside the union, and (3) gives
\[
                \alpha\text{ and }\gamma\text{ are closed on }S.
                                                               \tag{6}
\]

The boundary edge at $b_2$ has color $\beta$, whereas $\gamma$ is
missing there, so $\gamma\ne\beta$; also $\alpha\ne\beta$ is part of the
counterexample data.  Split according to $\alpha=\gamma$.  In the
equality branch put $\mu_1=\varphi$ and retain the original ETT,
hierarchy, maximum property, and (3) by reflexivity; no identity
``recoloring'' is passed to \noderef{OutsideKempeRecoloring}.  In the
other branch, the proved inequality $\alpha\ne\gamma$ is the first field
of \noderef{OutsideKempeRecoloring}; exchange these colors on all full
components disjoint from $S$.  Then
\noderef{OutsideKempeRecoloringStable} makes $\mu_1$ proper, leaves
exactly $e_0$ uncolored, fixes every exact edge incident with $S$, and is
stable on every displayed stage.  Applying
\noderef{HierarchyLEMStable} to the truncated ETT on $S$ checks the
lower package, maximum property, positivity, properness,
one-uncoloredness, hierarchy base, good hierarchy, and the recoloring.
Its output is a corresponding same-rung ETT under $\mu_1$, with maximum
property, identical stages and $\Gamma$-sets, and the transported
instance of (3).  Thus both equality branches provide exactly the same
named data.

Because both colors in (6) have empty boundary, every
$\alpha$-edge on the portion of $P_1$ or $P_2$ outside $S$ lies in a
full $(\alpha,\gamma)$ component disjoint from $S$.  Inspection of
\noderef{TwoColorComponentSwap} therefore changes precisely those
external $\alpha$-edges to $\gamma$ and fixes the $\beta$-edges: the
latter assertion uses both $\beta\ne\alpha$ and $\beta\ne\gamma$ proved
above.  In the equality branch the same paths are unchanged.  Hence
$P_1,P_2$ are two distinct $(\gamma,\beta)$ exit paths under $\mu_1$.

Let $K_2$ be the full $(\gamma,\beta)$ component through $b_2$ under
$\mu_1$, and let
$\mu_2=\mu_1/K_2$.  The vertex $b_2$ misses $\gamma$ and is incident
with the first $\beta$-edge of $P_2$, so properness makes its
incidence-graph degree one.  The outside end $a_2$ misses one of the two
path colors by \noderef{TreeExitPath}; because the path is nontrivial it
is incident with the other color, and it too has degree one.  Every node
has degree at most two by
\noderef{TwoColorIncidenceDegreeAtMostTwo}.  A node of $K_2$ outside the
simple path $P_2$ would attach at an internal node, producing degree
three, or at an endpoint, producing degree two there.  Both are
impossible.  Thus $P_2$ exhausts $K_2$.  Its exit-path clause says that
its only vertex in $S$ is $b_2$, so $K_2$ is disjoint from the prefix
strictly before $b_2$.  This swap is therefore an outside Kempe
recoloring for that shorter prefix.  The colors are distinct because one
is missing and the other is present on the first boundary identity at
$b_2$.  Therefore
\noderef{TwoColorComponentSwapProper} gives properness of $\mu_2$;
\noderef{TwoColorComponentSwap} fixes the value $\mathsf{none}$ and only
transposes two nonempty colors, so \noderef{IsOneEdgeUncolored} shows
that exactly $e_0$ remains uncolored.

Let $B$ be the prefix immediately before $b_2$.  If $A\subsetneq B$,
use the valid $p-1$-divider presentation ending at $B$ described above.
If $B=A$ and $p>1$, use instead the valid $p-2$-divider presentation
ending at $T_{n,p-1}=A$, apply \noderef{HierarchyLEMStable} there, and
only afterwards append the exact edge and vertex $b_2$.  If $B=A$ and
$p=1$, use the separate base route through
\noderef{OutsideKempeRecoloringStable}, \noderef{Td}, and
\noderef{TechnicalTheoremBase}.  In all three branches the one-edge
reconstruction, \noderef{Td}, and
\noderef{MaximumPropertyTransport} obtain a same-rung maximum ETT
$C=S(b_2)$ under $\mu_2$.  Adding $b_2$ makes the last stage-to-terminal
containment strict, so its retained hierarchy now has exactly $p-1$
dividers, the same stages through $A$, the same $\Gamma$-sets, the old
internal minus-closed fields, and the final avoidance field checked on
the exact new edge.  These are precisely all properness,
one-uncoloredness, corresponding-ETT, maximum-property, same-stage,
same-$\Gamma$, and minus-closed obligations; none is inferred merely
from the word ``stable.''

Now (5) is preserved on $A$, since the swapped component is disjoint
from the prefix before $b_2$.  The swap makes $\beta$ missing at $b_2$.
Both vertices lie in $C$ and are distinct.  Applying (6.1) to the newly
constructed good hierarchy, whose $p-1$ dividers are fewer than $q$,
makes $V(C)$ elementary under $\mu_2$.  The same color $\beta$ is
missing at the two displayed vertices, a contradiction.

\emph{Subcase 1.2: $b_2\in V(A)$.}
We first construct a color $\theta$ with the exact properties used in
the recoloring.  The stored-color bound
\noderef{HierarchyStoredColorsCardinality} gives
$|D_{n,j}|\leq n$ at every stage.  For each fixed $j$,
`twoColors' and `disjointPairs' in \noderef{GoodHierarchy} say that the
pairs indexed by $D_{n,j}$ have cardinality two and are pairwise
disjoint.  Hence
\[
 \left|\bigcup_{\eta\in D_{n,j}}\Gamma^j_\eta\right|
       =2|D_{n,j}|\leq2n.                                    \tag{7}
\]
This is the required use of the actual stored-pair fields, not an
uncited cardinality shortcut.

If $p=1$ and the terminal kind is not PE, (L) and (7) choose
\[
 \theta\in\overline\varphi(T_n)
       -\left(\bigcup_{\eta\in D_{n,1}}\Gamma^0_\eta
          \cup\{\alpha,\beta\}\right).
\]
Indeed the reserved union has at most $2n$ colors, so after excluding
$\alpha$ and $\beta$ at least nine of the $2n+11$ colors remain.
To verify closure on $T_n$, inspect the last non-PE constructor in
\noderef{TashkinovSeriesStep}.  In RE its generating coloring equals the
previous coloring, under which $T_n$ is explicitly closed.  In SE,
\noderef{MaximumDefectiveChoice} says that the selected generating
coloring is stable on $T_n$ relative to the previous coloring.  Its
vertex clause in \noderef{StableColoring} preserves the missing set;
properness and exact one-uncoloredness on both sides let the symmetry
clause of \noderef{sc2} transfer a hypothetical boundary identity of a
missing color back to the previous coloring, contradicting
\noderef{ColoredVertexSetClosed}.  Hence $T_n$ is closed under the actual
generating coloring in either non-PE case, while (3) closes $\theta$ on
$T_{n,1}$.
If $p=1$ and the terminal kind is PE, the `parallelBase' field of
\noderef{GoodHierarchy} gives at least four colors in
\[
 (\overline\varphi(T_n)\cap\overline\varphi(\mathcal B.\mathit{closure}))
       -\bigcup_{\eta\in D_{n,0}}\Gamma^0_\eta.
\]
Delete $\alpha$ and $\beta$ from this set and choose $\theta$ from the
at least two remaining colors.  It is missing on both $T_n$ and the base
closure and satisfies $\theta\ne\alpha,\beta$.  The pre-step
elementarity premise of \noderef{ParallelCommonMissingClosed} is lower
assertion (i) for the one-rung restriction, with
\noderef{TechnicalTheoremBase} at its zero-rung endpoint.  That lemma
therefore makes the singleton $\{\theta\}$ closed on the exact hierarchy
base core $T_n\vee R_n$.  Thus the missingness and singleton-closedness
hypotheses of the exact third conclusion of \noderef{interchange} are
both available.  Unfolding \noderef{TashkinovStoredColors} and
\noderef{HierarchyStoredColorsAt} and using
$T_{n,0}^{*}\subseteq T_{n,1}$ gives
$D_{n,1}\subseteq D_{n,0}$ directly; hence the chosen color avoids the
smaller terminal reserved union and clause (3) closes $\theta$ on
$T_{n,1}$.

If $p\geq2$, (L), inclusion of terminal-ladder missing colors along the
ordered stages, and (7) choose
\[
 \theta\in\overline\varphi(T_{n,p-2})
       -\left(\bigcup_{\eta\in D_{n,p-1}}\Gamma^{p-2}_\eta
          \cup\{\alpha,\beta\}\right).                       \tag{8}
\]
Again at least nine choices remain after the two explicit exclusions.
The minus-closed field at $p-1$ makes $\theta$ closed on
$T_{n,p-1}$.  For a stored color whose pair is not refreshed, the
`uniqueRefresh' field of \noderef{GoodHierarchy} identifies its new pair
with the old $\Gamma^{p-2}$-pair, which (8) avoids.  The unique refreshed
pair is missing on
$V(T_{n,p-1})-V(T_{n,p-2}^{*})$; if it contained $\theta$, then
$\theta$ would be missing both on $T_{n,p-2}$ and at a distinct new
vertex, contradicting elementarity of $T_{n,p-1}$.  Thus $\theta$ is not in
$\bigcup_{\eta\in D_{n,p}}\Gamma^{p-1}_\eta$; it is still missing on
$S$, and (3) closes it on $S$.  In all three alternatives we have
\[
 \begin{cases}
 \theta\in\overline\varphi(T_n),\quad
 \theta\text{ closed on }T_{n,0}^{*}\text{ and }T_{n,1},&p=1,\\
 \theta\in\overline\varphi(T_{n,p-1}),\quad
 \theta\text{ closed on }T_{n,p-1}\text{ and }S,&p>1.
 \end{cases}                                                  \tag{9}
\]

Our choice gives $\alpha\ne\theta$ and $\theta\ne\beta$.  Thus the first
inequality required by \noderef{OutsideKempeRecoloring} is proved before
we let $\mu_3$ exchange $\alpha,\theta$ on every full component outside
$S$.  By (6), \noderef{OutsideKempeRecoloringStable} makes $\mu_3$
proper, leaves exactly $e_0$ uncolored, and fixes all stage-incident edge
identities and missing sets.  The inequality $\theta\ne\beta$ and the
given $\alpha\ne\beta$ show directly from
\noderef{TwoColorComponentSwap} that every $\beta$-edge is fixed.
\noderef{HierarchyLEMStable}, applied with every hypothesis listed in
Subcase 1.1, supplies a corresponding same-rung maximum ETT and the same
good hierarchy, stored sets, terminal minus-closed certificate, and
$\Gamma$ data under $\mu_3$.

It remains to prove component failure rather than infer it merely from
two disjoint tails.  On $G-V(S)$ the swap gives an incidence-graph
isomorphism from the old $(\alpha,\beta)$ subgraph to the new
$(\theta,\beta)$ subgraph: old $\alpha$ edge-nodes become $\theta$
edge-nodes, $\beta$ edge-nodes and every vertex-node are fixed, and old
$\theta$ edge-nodes leave the two-color graph.  This follows edge by edge
from \noderef{TwoColorComponentSwap} and
\noderef{TwoColorIncidenceGraph}.

If $P_1,P_2$ came from distinct old full components, follow any proposed
new connecting path between their outside portions and choose its first
and last contacts with $S$.  Between those contacts it uses only the
unchanged internal $(\theta,\beta)$ incidence graph, while outside it
pulls back under the preceding isomorphism to the two disjoint old
components.  Both contacts are entered by distinct $\beta$ boundary
identities and $\theta$ has no boundary identity.  Deleting the internal
contact segment would therefore make the two pulled-back components
share a node, contrary to \noderef{TwoColorComponentsDisjoint}; if the
segment is nonempty, its two end incidences instead give two ends in
$A$ missing the same one of $\theta,\beta$, contrary to the elementarity
of $A$.  Thus the two transformed tails remain in distinct full
components.

If $P_1,P_2$ are the two tails of one old component, inspect its simple
spanning path between $b_1$ and $b_2$.  Both boundary identities have
color $\beta$, so alternation supplies an $\alpha$-edge identity on the
segment through $S$ immediately after one of the contacts.  That exact
edge is incident with $S$ and hence is fixed as $\alpha$ by the outside
recoloring; it is absent from the new $(\theta,\beta)$ incidence graph.
The old full component is a path and every node has degree at most two,
so there is no second route around this removed edge.  It separates the
two transformed outside tails into two full
$(\theta,\beta)$ components.  In both original-component cases the
transformed paths therefore witness genuine failure of
$(\theta,\beta)$ interchangeability on the preceding tree.

When $p=1$ and the terminal kind is PE, apply conclusion (iii) of
\noderef{interchange} to the truncated ETT: properness and exact
one-uncoloredness are supplied by the preceding outside-recoloring
lemma; \noderef{HierarchyOutsideKempeStable} supplies the required
combined stability on the exact terminal $T_n$ and base closure; the
terminal-PE equation, hierarchy base, maximum property, positivity, and
the complete lower six-assertion package are the stated hypotheses; and
(9) supplies missingness and singleton closedness on `base.core'.  This
contradicts the two paths.  When $p=1$ and the kind is not PE, the
terminal ladder is the starred base.  If $n>1$, restrict the certified
series by one rung and display $T_n$ as its final ETT.  Its maximum
property is the restriction of \noderef{MaximumProperty}.  In RE the
generating color is already the preceding color; in SE the
`selectedStable' field of \noderef{MaximumDefectiveChoice}, followed by
transitivity in \noderef{sc2}, combines with the outside stability of
$\mu_3$ to give exactly the preceding-rung stability relation.  Lower
assertion (ii) therefore applies to $\mu_3$.  If $n=1$, the same
interchangeability conclusion, with its properness, one-uncoloredness,
and stability premises just checked, is
\noderef{TechnicalTheoremBase}.  In either event the two paths contradict
interchangeability.  When $p>1$, truncate the transported hierarchy at
$T_{n,p-1}=A$.  Its divider count is $p-2$, its terminal
minus-closed certificate is the old field at $p-2$, (9) supplies
$\theta$-missingness and singleton closedness there, and the large bound
is fixed because every edge incident with $T_n$ was fixed.  It is
therefore an admissible presentation of strictly smaller primary measure,
contradicting the lexicographic minimum.  This
finishes Case 1.

\textbf{Case 2: $\beta\notin\overline\varphi(S)$.}
No vertex of $S$ misses $\beta$.  Counting $\beta$-incidences now gives
\[
                  |\partial_{\varphi,\beta}(S)|\equiv
                  |V(S)|\equiv1\pmod2.                      \tag{10}
\]
No nontrivial path component meeting $S$ can be wholly inside it: such a
path has two endpoints missing one of its two colors, whereas $\beta$ is
missing nowhere on $S$ and elementarity permits $\alpha$ to be missing
at only one vertex.  Thus the two witness components supply crossings.
Equation (10) raises their total number to at least three.  In each
max-degree-two path component take the portions from its outside
endpoints to its first encountered vertices of $S$.  Portions in
different full components are disjoint by
\noderef{TwoColorComponentsDisjoint}.  The two end portions of one full
path component are disjoint as well: if their first contacts coincided,
two distinct $\beta$ boundary identities would meet that vertex, since
$\alpha$ has no boundary edge, contradicting properness.  Thus these
outer portions are pairwise vertex-disjoint.  The odd total and the two
witness components provide at least three of them.  Select three and call them
$P_1,P_2,P_3$, with outside ends $a_j$, inside ends $b_j$, and distinct
exact first boundary-edge identities $f_j$, ordered so that
$b_1\prec b_2\prec b_3$.  The paths, their outside ends, and their first
edge identities are pairwise distinct, though two paths may arise from
the two ends of one full component.

The secondary coordinate of the global comparison was defined precisely
from such triples.  Replace the three just selected, if necessary, by a
triple with least total length.  Thus
$|P_1|+|P_2|+|P_3|$ realizes the secondary coordinate minus one.  No
coloring or hierarchy is held fixed in this minimum: another proper
one-uncolored generating coloring, arbitrary certified input,
corresponding same-rung maximum ETT, good hierarchy with the same
primary divider count, terminal minus-closed certificate, lower package,
(6.1), and large bound belongs to the same comparison class.

\emph{Subcase 2.1: $b_3\notin V(A)$.}
Choose $\gamma\in\overline\varphi(b_3)$ by the same certified degree
argument as in Subcase 1.1.  Elementarity and
\noderef{GoodHierarchyReservedDisjoint} again imply that $\gamma$ is
outside the reserved union in (3), so (3) closes it on $S$.  Since the
first boundary identity at $b_3$ has color $\beta$ while $\gamma$ is
missing there, $\gamma\ne\beta$.  Split $\alpha=\gamma$ exactly as in
Subcase 1.1.  In the equality branch put $\mu_4=\varphi$ and use the
identity ETT and hierarchy data.  In the inequality branch first record
$\alpha\ne\gamma$, then recolor all outside
$(\alpha,\gamma)$ components.  The exact application of
\noderef{HierarchyLEMStable} is the one already checked for $\mu_1$: it
gives properness, exact one-uncoloredness, a corresponding same-rung ETT,
maximum property, the same stages and $\Gamma$-sets, and the transported
terminal minus-closed data.  In both branches the three $P_j$ become
$(\gamma,\beta)$ exit paths.  Their $\beta$-edges are fixed because
$\beta$ differs from both $\alpha$ and $\gamma$.

Swap the full $(\gamma,\beta)$ component through $b_3$ and call the
result $\mu_5$.  As at $b_2$, $b_3$ has incidence degree one because it
misses $\gamma$ and is incident with a $\beta$-edge, while the outside
end of $P_3$ has degree one by the exit-path condition.  The bound
\noderef{TwoColorIncidenceDegreeAtMostTwo} shows that $P_3$ exhausts
this component.  Hence the component meets $S$ only at $b_3$ and is
disjoint from the prefix $B$ before $b_3$.  If $A\subsetneq B$, use the
valid $p-1$-divider presentation ending at $B$.  If $B=A$ and $p>1$,
use the $p-2$-divider presentation ending at $A$ before restoring the
next edge; if $B=A$ and $p=1$, use the separate base route.  Thus
\noderef{TwoColorComponentSwapProper},
\noderef{IsOneEdgeUncolored}, \noderef{HierarchyLEMStable},
\noderef{Td}, and \noderef{MaximumPropertyTransport} construct the exact
prefix $C=S(b_3)$ as a corresponding same-rung maximum ETT with the
same stages through $A$, the same $\Gamma$ data, all retained
minus-closed fields, strict final containment, and the final avoidance
check on the reconstructed edge.  The swap makes $\beta$ missing at
$b_3$.

Adjoin the two individual edge identities $f_1$ and $f_2$, in that
order, to $C$.  Their old endpoints are $b_1,b_2\in V(C)$, their other
endpoints are the next vertices on the two simple exit paths, and they
are distinct even when two tails came from one full component: the two
end portions of its simple spanning path are disjoint.  Hence each
application of \noderef{TreeSequenceAdjoinEdge} retains the ordered tree
and does not identify parallel edges.  Call the resulting tree $J$.

Before any iteration, construct the data on $J$ in dependency order.
The swap component is disjoint from $B\supseteq T_n$, so it is an
outside $(\gamma,\beta)$ recoloring relative to $T_n$; the inequality
$\gamma\ne\beta$ was proved before the swap.  Hence
\noderef{OutsideKempeRecoloringStable} gives terminal-ladder stability
of $\mu_5$, and \noderef{Td} gives a \noderef{ModColoring} of the
original same-rung ETT.  Compose the exact augmentation sequence ending
at $C$ with the two \noderef{TashkinovAugmentation} steps for
$f_1,f_2$.  The universal conclusion in \noderef{ModColoring} first
supplies the \noderef{CorrespondingExtendedTashkinovTree} on $J$;
\noderef{MaximumPropertyTransport} then supplies maximum property.

Now copy the stages through $A$ and take $J$ as terminal.  This gives a
\noderef{TashkinovHierarchy} on that replacement ETT with $p-1$
dividers: base stage, divider positions, nestedness, and all old strict
containments are unchanged, while $A\subsetneq C\subsetneq J$ proves the
final strictness.  The `ladderTreesAgree', `storedColorsAgree',
`connectingEdgesAgree', and `extensionKindsAgree' fields of
\noderef{CorrespondingExtendedTashkinovTree}, followed by the equalities
in \noderef{SameHierarchyStages} and \noderef{SameGoodHierarchy},
identify the replacement raw stored sets, stages, hierarchy stored sets,
and paper-relevant $\Gamma$ values with the transported ones.  Thus pair
sizes, pairwise disjointness, unique refresh, parallel base, and every
internal minus-closed clause are the transported old fields.

The final avoidance field is not simply copied.  The old avoidance field
controls the segment from $A$ through $C$, because the recoloring fixes
all displayed edges there and all first-missing positions for stored
colors.  Both new exact identities $f_1,f_2$ have color in
$\{\beta,\gamma\}$.  The color $\gamma$ was proved outside every
$\Gamma^{p-1}_\eta$; if $\beta$ belonged to one of these pairs,
\noderef{GoodHierarchyReservedDisjoint} would make it missing on $A$
and hence on $S$, contrary to Case 2.  Therefore neither new edge color
belongs to any indexed pair, so unfolding
\noderef{HierarchySegmentColors} verifies avoidance on $J$.  This proves
all five fields of \noderef{GoodHierarchy} for $J$ before the iteration
starts.  Corresponding ETT and maximum property were proved first, so
the construction respects the dependency order.

Starting from $J$, perform the manuscript's finite Step-2 augmentation.
At a current tree $U$, append an exact boundary edge $f$ only when its
$\mu_5$-color is missing on $U$ and, for every
$\eta\in D_{n,p-1}$, the appended segment before the first vertex
missing $\eta$ uses no color in $\Gamma^{p-1}_\eta$.  This condition is
exactly the avoidance predicate encoded by
\noderef{HierarchySegmentColors}.  Each step is a
\noderef{TashkinovAugmentation}: the old endpoint is in $U$, the other
endpoint is new, the exact edge identity is appended, and its displayed
color has a missing witness in $U$.  Since $E$ and $V$ are finite and
each step strictly increases the ordered vertex list, a maximal finite
sequence exists.  Induct along it starting with the data on $J$ just
verified.  At a step the accumulated exact edges first form a
\noderef{TashkinovAugmentationSequence}; the universal clause of
\noderef{ModColoring} supplied by \noderef{Td} gives the corresponding
ETT at the new tree, and \noderef{MaximumPropertyTransport} gives
maximum property.  On that ETT copy the base stage, dividers, and
internal stages to construct its hierarchy.  Pair sizes, pairwise
disjointness, unique refresh, parallel base, and internal
minus-closedness are unchanged.  The selection predicate is precisely
the new terminal avoidance field after unfolding
\noderef{HierarchySegmentColors}.  Hence all five fields of
\noderef{GoodHierarchy} hold after the append.  Let $T'$ be
a maximal last tree.  Its corresponding ETT, maximum property,
hierarchy, same stages and $\Gamma$-sets, and all five good-hierarchy
fields have therefore been proved before we invoke terminal closure or
(6.1).

Put
\[
 D'_{n,p}=\left(\bigcup_{h\leq n}S_h\right)
               -\overline{\mu_5}(T').
\]
This is exactly the terminal stored-color set of the hierarchy just
constructed.  Indeed, unfold \noderef{TashkinovStoredColors}: at the
positive last rung it is the cumulative union of the exact raw stored
sets $S_h$ through $n$, with the colors missing on $T_n$ deleted.  The
corresponding-ETT fields preserve each raw stored set, and $\mu_5$ fixes
every edge incident with $T_n$.  Next unfold
\noderef{HierarchyStoredColorsAt}: at the preceding starred stage it is
this ladder stored set with $\overline{\mu_5}(A)$ deleted.  Since
$T_n\subseteq A\subseteq T'$, missing-color monotonicity and elementary
set difference give
\[
 D'_{n,p}=D_n^{\mu_5}-\overline{\mu_5}(T')
       \subseteq D_n^{\mu_5}-\overline{\mu_5}(A)=D_{n,p-1}.
\]
Only now, after the hierarchy exists, may
\noderef{HierarchyStoredColorsAntitone} also be applied to its internal
stages; the displayed terminal inclusion itself follows from the two
unfolded definitions.

If a color missing on
$T'$ but outside
$\bigcup_{\eta\in D'_{n,p}}\Gamma^{p-1}_\eta$ had a colored boundary
identity of color $c$, that identity would satisfy the Step-2 rule.
Indeed, fix $\eta\in D_{n,p-1}$.  If $\eta\notin D'_{n,p}$, then
$\eta$ is already missing on $T'$, so its first missing vertex precedes
the proposed new last edge and that edge does not enter the tested
segment.  If $\eta\in D'_{n,p}$, the hypothesis on $c$ says
$c\notin\Gamma^{p-1}_\eta$, while all older segment edges already pass
the avoidance test.  Thus every $\eta$ passes, and maximality forbids
the boundary identity.  This gives the exact statement
\[
 T'\text{ is }
 \left(\bigcup_{\eta\in D'_{n,p}}\Gamma^{p-1}_\eta\right)^{-}
 \text{-closed under }\mu_5.                                 \tag{11}
\]
The hierarchy, all five good-hierarchy fields, corresponding ETT, and
maximum property were established before defining $D'_{n,p}$; (11) is
the additional terminal certificate (C).  Its $p-1<q$ dividers now let
(6.1) conclude
\[
                       V(T')\text{ is elementary under }\mu_5. \tag{12}
\]

We show that no $a_j$ lies in $T'$.  If some $a_j$ did, the endpoint
property of the path component says that it misses $\gamma$ or $\beta$.
It cannot miss $\beta$, because $b_3$ already misses $\beta$ and (12)
is elementary; hence it misses $\gamma$.  Both colors are outside the
reserved pairs: $\gamma$ was excluded by elementarity; if $\beta$ lay in
a $\Gamma^{p-1}$-pair, \noderef{GoodHierarchyReservedDisjoint} would
make it missing on the preceding starred stage and hence on $S$,
contrary to Case 2.  Since $D'_{n,p}\subseteq D_{n,p-1}$, both colors
are outside the union occurring in (11).  Equation (11), the
$\gamma$-missing witness $a_j$, and the $\beta$-missing witness $b_3$
therefore make both colors closed on $T'$.

Each $P_i$ meets $T'$ at $b_i\in S$.  If its other end were outside
$T'$, the first edge on leaving $T'$ would be a $\gamma$- or
$\beta$-colored boundary identity, contradicting the two closedness
conclusions.  Hence every $a_i$ lies in $T'$.  No $a_i$ can miss
$\beta$, because $b_3$ already does and (12) is elementary; consequently
all three distinct outside ends miss $\gamma$.  Already two of them
contradict (12).  Thus no $a_j$ lies in $T'$.

For each $j$, take the last portion $L_j$ of $P_j$ from its last vertex
in $T'$ to $a_j$.  It is a $T'$-exit path.  The three paths have distinct
outside endpoints, while a finite max-degree-two path component has only
two endpoints; hence at least two full components meet $T'$, giving the
failure of interchangeability required in the comparison class.  The
identity $f_1$ was deliberately
adjoined to $J\subseteq T'$, so it is not in $L_1$; all other deletions
are nonnegative.  Consequently
\[
       |L_1|+|L_2|+|L_3|<|P_1|+|P_2|+|P_3|.                 \tag{13}
\]
The first comparison color must now be $\beta$, not $\gamma$.  The swap
makes $\beta$ missing at $b_3\in V(T')$.  We proved above that $\beta$
belongs to no $\Gamma^{p-1}_\eta$, hence in particular to none indexed
by $D'_{n,p}$.  Unfolding \noderef{TreeColorSetMinusClosed} in (11)
therefore proves that the singleton $\{\beta\}$ is closed on $T'$.
Also $\beta\ne\gamma$, as was proved before the component swap.

The $(\gamma,\beta)$ and $(\beta,\gamma)$ incidence graphs are literally
equal: unfold \noderef{TwoColorIncidenceGraph} and commute the two color
disjuncts.  The definitions \noderef{TwoColorComponent},
\noderef{TwoColorPathComponent}, and \noderef{TreeInterchangeable} are
therefore symmetric in the two colors.  Consequently the at least two
full components containing the $L_j$ prove failure of
$(\beta,\gamma)$ interchangeability.

We verify every remaining comparison-class field.  Properness and exact
one-uncoloredness follow from \noderef{TwoColorComponentSwapProper} and
\noderef{IsOneEdgeUncolored}.  The corresponding same-rung ETT,
maximum property, good hierarchy, $p-1$ divider count, same stages,
stored sets and $\Gamma$ family, and all five good fields were proved in
dependency order before (11).  Equation (11) is terminal certificate
(C), and (12) records the permitted use of (6.1).  The lower package is
the fixed universal hypothesis.  Both recolorings fix every edge
incident with $T_n$: the first is outside $S$, and the second component
is disjoint from $B\supseteq T_n$.  Thus the large bound is unchanged.
Finally the $L_j$ are three distinct exit paths with distinct outside
ends in at least two full components.

Hence $(\mu_5,T',\beta,\gamma,L_1,L_2,L_3)$ is an admissible
presentation of the same primary measure.  If $\gamma$ is missing on
$T'$, its secondary measure is zero, strictly below the positive Case 2
measure.  If $\gamma$ is not missing, (13) makes its secondary measure
strictly smaller.  Either alternative contradicts the global
lexicographic minimum.

\emph{Subcase 2.2: $b_3\in V(A)$.}
We repeat every choice.  At a fixed stage $j$, the theorem
\noderef{HierarchyStoredColorsCardinality} gives $|D_{n,j}|\leq n$.
The `twoColors' and `disjointPairs' fields of
\noderef{GoodHierarchy} make the union of the pairs indexed by this set
have cardinality exactly $2|D_{n,j}|$, hence at most $2n$.

Suppose first that $p=1$ and the terminal kind is not PE.  The
$2n+11$ colors in (L) therefore include a color
\[
 \theta\in\overline\varphi(T_n)-
      \left(\bigcup_{\eta\in D_{n,1}}\Gamma^0_\eta
        \cup\{\alpha,\beta\}\right).                         \tag{14}
\]
The two extra exclusions still leave at least nine choices.
In RE the terminal generating coloring is the previous coloring and the
constructor of \noderef{TashkinovSeriesStep} states that $T_n$ was
closed.  In SE, the `selectedStable' field of
\noderef{MaximumDefectiveChoice}, symmetry from \noderef{sc2}, and the
old closedness transfer every hypothetical boundary identity of a
missing color back to the previous coloring.  Thus $\theta$ is closed on
$T_n=T_{n,0}^{*}$.  Equation (14) and the minus-closed statement (3)
also make it closed on $S=T_{n,1}$.

Suppose next that $p=1$ and the terminal kind is PE.  The `parallelBase'
field of \noderef{GoodHierarchy} supplies at least four colors in
\[
 (\overline\varphi(T_n)\cap
   \overline\varphi(\mathcal B.\mathit{closure}))-
      \bigcup_{\eta\in D_{n,0}}\Gamma^0_\eta.
\]
Choose $\theta$ after deleting $\alpha$ and $\beta$; the cardinality
four lower bound leaves at least two choices.  The
pre-step elementarity required by \noderef{ParallelCommonMissingClosed}
is lower assertion (i) for $n>1$ and
\noderef{TechnicalTheoremBase} for $n=1$.  Hence $\theta$ is missing and
singleton-closed on the exact base core $T_n\vee R_n=T_{n,0}^{*}$.
Unfolding \noderef{TashkinovStoredColors} and
\noderef{HierarchyStoredColorsAt} and using
$T_{n,0}^{*}\subseteq T_{n,1}$ gives
$D_{n,1}\subseteq D_{n,0}$ directly.  The chosen color is therefore
outside the reserved union in (3), so it is also closed on
$S=T_{n,1}$.

Finally suppose $p>1$.  Missing-color monotonicity along the ordered
stages and (L) give at least $2n+11$ colors missing on $T_{n,p-2}$.
Delete the at most $2n$ colors in
$\bigcup_{\eta\in D_{n,p-1}}\Gamma^{p-2}_\eta$ and also delete
$\alpha,\beta$.  At least nine colors remain; choose $\theta$ there.
The minus-closed field at $p-1$ closes $\theta$ on
$T_{n,p-1}$.  For every nonrefreshed stored color, `uniqueRefresh'
identifies its $\Gamma^{p-1}$-pair with the old pair just avoided.  The
refreshed pair is missing on the new segment above
$T_{n,p-2}^{*}$; membership of $\theta$ in it would give one
$\theta$-missing vertex in that segment and another in
$T_{n,p-2}$, contradicting elementarity of $T_{n,p-1}$.  Thus $\theta$
lies outside the reserved union in (3), remains missing on $S$, and is
closed there.  We have proved, independently in this subcase, the two
alternatives displayed in (9).

Now $\alpha\ne\theta$ is explicit, so exchange $\alpha,\theta$ on every
full component outside $S$ and call the coloring $\nu$.  Both colors are
closed on $S$, and the proved inequality supplies the first field of
\noderef{OutsideKempeRecoloring}.  The theorem
\noderef{OutsideKempeRecoloringStable} proves that $\nu$ is proper,
leaves exactly $e_0$ uncolored, fixes every stage-incident edge identity
and every stage missing set, and is stable for every required control
set.  Applying \noderef{HierarchyLEMStable} to the exact truncated ETT
on $S$ supplies a corresponding same-rung ETT under $\nu$, maximum
property, the identical ordered hierarchy and $\Gamma$ family, and the
transported minus-closed field.  Since the three $P_j$ have their inside
ends in $A$, all their external $\alpha$-edges lie in components disjoint
from $S$ and become $\theta$-edges.  Their $\beta$-edges and exact inside
portions are fixed because $\theta\ne\beta$ and $\alpha\ne\beta$.
Therefore the same three pairwise disjoint exit paths have three distinct
outside endpoints.  A full path component has only two end vertices by
\noderef{TwoColorComponentSpanningPath} and
\noderef{TwoColorIncidenceDegreeAtMostTwo}; hence the three outside ends
lie in at least two full $(\theta,\beta)$ components.  Thus
interchangeability fails on the preceding tree.

In the PE case with $p=1$, apply conclusion (iii) of
\noderef{interchange} to the exact base core.  Its terminal-PE equation,
hierarchy base, maximum property, positivity, and complete lower package
are the hypotheses of this lemma.  Properness and exact
one-uncoloredness of $\nu$ were just proved, and
\noderef{HierarchyOutsideKempeStable} gives the required combined
stability on the exact terminal ladder and base closure.  The PE choice
above, transported across the incident-edge and missing-set equalities,
gives missingness and singleton closedness on `base.core'.  The resulting
interchangeability contradicts the two components.

In the non-PE case with $p=1$, display $T_n$ as the terminal tree of the
one-rung restriction when $n>1$.  Maximum property restricts, and the RE
color equality or the SE `selectedStable' field combined with
transitivity in \noderef{sc2} gives precisely the preceding-rung stable
relation for $\nu$.  Lower assertion (ii) makes $\theta,\beta$
interchangeable on $T_n$; when $n=1$, this is exactly the second
conclusion of \noderef{TechnicalTheoremBase}.  Again the two components
contradict it.

If $p>1$, truncate the transported hierarchy at $T_{n,p-1}$.  It has
$p-2$ dividers, the old terminal minus-closed field, the missingness and
closedness of $\theta$, the same lower package and (6.1), and the same
large terminal-ladder bound because the outside recoloring fixes every
edge incident with $T_n$.  The two components therefore give an
admissible counterexample of smaller primary measure, contradicting the
global lexicographic minimum.  This disposes of
Subcase 2.2 without appealing to omitted manuscript details.

All possibilities following the assumed failure have produced a
contradiction.  Therefore $\alpha$ and $\beta$ are
$T_{n,p}$-interchangeable under $\varphi_n$, as claimed.
\end{proof}
